% !TEX root = ../linnik399.tex
\section{Far density and representatives}\label{sec:far}

The per-character costs must be summed over many characters. A weighted density
estimate supplies a common budget for that sum. Its zeros have physical height at most
$1$, so we also specify how to choose representatives within this height range.

\subsection{Xylouris's weighted density lemma}
The following is Xylouris's improvement \cite[Lemma~5.1, (5.19), p.~66]{Xylouris2011nullstellen} of
Heath-Brown's Lemma~11.1 \cite{HeathBrown1992zero}; the latter is the case $J=1$, $w_0\equiv1$.
(Xylouris writes $M$ for our $J$; we reserve $M$ for a separation constant.)

\begin{inp}[Far density]\label{inp:far}
Let $\eps,c_1,c_2>0$, $J\in\NN$, and $\alpha_1,\dots,\alpha_J\ge0$ with $\sum_i\alpha_i=1$. Put
\[
  x=\tfrac23+3c_1+c_2,\qquad u_i=\tfrac13+2c_1+\frac{ic_2}{J}\quad(0\le i\le J),
\]
and let $w_0:[u_0,x]\to\RR$ be continuous and continuously differentiable except at finitely many
points, with $1\ll w_0\ll1$ and $w_0'\ll1$. Choose for each nonprincipal character $\chi$ modulo
$q$ at most one zero $\rho_\chi$ of $L(s,\chi)$ with $|\Im\rho_\chi|\le1$ and
$\lambda_\chi\le\lambda_0=\frac13\log\log\LL$. Then for $q\ge q_0$, with $q_0$ depending on all
the parameters,
\[
  \sum_\chi\Bigl(\int_{u_0}^xw_0(t)^2e^{2\lambda_\chi t}\,dt\Bigr)^{-1}\le
  \frac{J^2+\eps}{c_1c_2^2}\sum_{i=1}^J\alpha_i^2\int_{u_{i-1}}^xw_0(t)^{-2}
  \min\{t-u_{i-1},\,u_i-u_{i-1}\}\,dt .
\]
\end{inp}

In \cite{Xylouris2011nullstellen} the statement is made for the zeros $\rho^{(k)}$ that are chosen
for the characters counted by $N(\lambda_0)$, one zero per character \cite[\S3.2.2,
pp.~25--26]{Xylouris2011nullstellen}, \cite[\S11]{HeathBrown1992zero}; any such choice is allowed, and
the estimates in the proof \cite[pp.~67--72]{Xylouris2011nullstellen} are uniform over
these choices. This uniformity is needed because the selected zeros, including the
$T^*(q)$-representatives below, depend on $q$. A nonreal character and
its conjugate are different characters, and each contributes its own term; the principal
character is excluded, since $L(s,\chi_0)$ has no zero in the relevant region for large $q$
\cite[p.~21]{Xylouris2011nullstellen}. The hypotheses on $w_0$ are satisfied by our profiles, which
are those of Xylouris's own application \cite[(6.20)--(6.21), p.~80]{Xylouris2011nullstellen}.

\subsection{The two weight profiles}
We use two profiles of the shape \cite[(6.20)]{Xylouris2011nullstellen}, both with $J=10$ and
$\eps_0=10^{-7}$:
\[
  w_0(t)^2=e^{-\theta_{\mathrm{far}}t}\sqrt{\min(t-u_0+\eps_0,\ c_2+\eps_0)}\qquad(u_0\le t\le x).
\]
\emph{Profile~I} has $c_1=0.09035$, $c_2=0.235968$ and $\theta_{\mathrm{far}}=1.28683$, and
\emph{profile~II} has $c_1=0.0821922$, $c_2=0.2170903$ and $\theta_{\mathrm{far}}=1.4964274$. The weights
$\alpha_i$ are the exact decimals of \cref{tab:alpha}; in each profile they sum to $1$. Both
profiles satisfy the hypotheses of \cref{inp:far}: $w_0$ is continuous and positive on
$[u_0,x]$, and it is smooth except at $t=u_0+c_2$.

\begin{table}[ht]
\centering
\small
\begin{tabular}{rll}
\toprule
$i$ & profile~I $\alpha_i$ & profile~II $\alpha_i$ \\
\midrule
1 & 0.0788827218 & 0.0806195583 \\
2 & 0.0849386148 & 0.0862705300 \\
3 & 0.0895629779 & 0.0905489307 \\
4 & 0.0938231516 & 0.0944644867 \\
5 & 0.0979710491 & 0.0982543287 \\
6 & 0.1021284916 & 0.1020315591 \\
7 & 0.1063732939 & 0.1058668825 \\
8 & 0.1107651869 & 0.1098131140 \\
9 & 0.1153565492 & 0.1139152197 \\
10 & 0.1201979632 & 0.1182153903 \\
\bottomrule
\end{tabular}
\caption{The weights $\alpha_i$ of the two far profiles.}
\label{tab:alpha}
\end{table}

Put
\[
  w(\lambda)=\Bigl(\int_{u_0}^xw_0(t)^2e^{2\lambda t}\,dt\Bigr)^{-1},
\]
a positive decreasing function of $\lambda$. Let $V$ be the right-hand side of \cref{inp:far}
with $\eps=0$. Rigorous enclosures give
\[
  V=175.26640331\ldots\ \text{(profile I)},\qquad V=243.32098105\ldots\ \text{(profile II)}.
\]
Taking $\eps=\eta J^2$ in \cref{inp:far}, we obtain for $q\ge q_0$ and every choice of zeros as in
\cref{inp:far}:
\begin{equation}\label{eq:farbudget}
  \sum_\chi w(\lambda_\chi)\le(1+\eta)V .
\end{equation}
Each leaf of the case tree uses one of the two profiles, recorded in its data. Both
are fixed finite choices within \cref{inp:far}.

\begin{lemma}[Comparison of the prime weight and the far weight]\label{lem:decay}
For both profiles, $x\le1.1737$. Hence $A>2x$ for $L=3.99$, and for every $\lambda\ge0$ the function
$\lambda\mapsto e^{-A\lambda}/w(\lambda)$ is decreasing. In particular
$e^{-A\lambda}\le w(\lambda)/w(0)$.
\end{lemma}

\begin{proof}
The values are $x=1.1736846\ldots$ and $x=1.1303335\ldots$, so $2x<2.348<A=3.156342$. Now
$e^{-A\lambda}/w(\lambda)=\int_{u_0}^xw_0(t)^2e^{(2t-A)\lambda}\,dt$, and $2t-A<0$ on
$[u_0,x]$.
\end{proof}

With $w(0)=10.7054\ldots$ (profile~I) and $w(0)=13.1730\ldots$ (profile~II), \cref{lem:decay} and
\eqref{eq:farbudget} give, for any selection satisfying \cref{inp:far},
\begin{equation}\label{eq:Kfar}
  \sum_\chi e^{-A\lambda_\chi}\le\frac{(1+\eta)V}{w(0)}=K_{\mathrm{far}},\qquad
  K_{\mathrm{far}}<16.38\ \text{(profile I)},\quad K_{\mathrm{far}}<18.48\ \text{(profile II)}.
\end{equation}
This bound controls errors that are summed over unboundedly many characters.

\subsection{Representatives}\label{subsec:representatives}
The leaf programs describe an ordinary character through one parameter, the parameter of a
\emph{representative} zero. For the near-density argument of \cref{sec:near} the representative
must be separated in height from any zero with a smaller parameter, and we arrange this by
choosing the height of the strip in which representatives are taken.

Fix $s_{\max}=3$; every anchor used in the near rows is at most $2.3$ (\cref{sec:rows}), and every
ordinary entry uses an anchor at most $s_{\max}$. Let $K_0$ be an upper bound for
the number of zeros, of all nonprincipal characters modulo $q$ and counted with multiplicity,
with $\lambda\le s_{\max}$ and $|\gamma|\le2$. By a log-free zero-density estimate
(\cref{inp:logfree}), $K_0$ may be taken independent of $q$. Next fix $M$ as in \cref{sec:join};
it depends on $K_0$.

\begin{lemma}[A strip boundary separated from low-parameter zeros]\label{lem:Tstar}
Let $K_0$ bound the number of zeros with $\lambda\le s_{\max}=3$ and $|\gamma|\le2$,
as above, and fix $M>0$. If $\LL>4M(K_0+1)$, there is $T^*\in[1/2,1]$ such that no nonprincipal $L(s,\chi)$ has a zero with
$\lambda\le s_{\max}$ and $\bigl||\mu|-T^*\LL\bigr|<M$. We take $T^*=T^*(q)$ to be the least such number.
\end{lemma}

\begin{proof}
The at most $K_0$ zeros in question exclude at most $K_0$ open intervals of $T$ of length $2M/\LL$
each. The interval $[1/2,1]$ has length $1/2>2MK_0/\LL$, so the remaining set is a nonempty
finite union of closed intervals, and it has a least element.
\end{proof}

\begin{definition}[Representatives]\label{def:representative}
For a nonprincipal character $\chi$, its \emph{$T^*$-representative} is a zero of $L(s,\chi)$ in
$R(l)$ with $|\gamma|\le T^*$ and least parameter, if there is one; its parameter is $\lambda_\chi$.
Representatives are used for the characters outside the first family; in \cref{thm:large} all
characters use their height-one representatives. The \emph{height-one representative} is defined in the same way with the bound
$|\gamma|\le1$. The reserved family of a leaf (\cref{subsec:vocabulary,def:specification}) and an inside first family use
height-one representatives.
\end{definition}

Every representative has $|\gamma|\le1$, so the far bound \eqref{eq:farbudget} applies to any
selection of representatives, one per character. Representatives for conjugate characters
may be chosen as conjugate zeros, with the same parameter, because every selection
region is symmetric in height.
The $T^*$-representative of $\chi$ satisfies $\lambda_\chi\ge$ the parameter of its height-one
representative. Shrinking the strip can only increase the minimum parameter; in terms
of real parts, the selected zero can only move to the left.

\subsection{The far constraint of a leaf}
In a leaf program (\cref{subsec:vocabulary}), the characters other than the first family are placed
in bins by the parameters of their representatives, and each bin has a column
(\cref{sec:lp,sec:casetree}). The column of a bin $[\lo,\hi]$ has far weight
$W_c=\lfloor Sw(\hi)\rfloor$; since $w$ is decreasing, each character in the bin contributes at least
$w(\hi)$ to the left side of \eqref{eq:farbudget}. The
\emph{tail column} collects the ordinary characters $\chi$ with a zero in $R_P$ and $\lambda_\chi\ge R$, where
$R=\max(3,r)$ and $r$ is the leaf's lower bound for ordinary representatives. Its value is the mass
$x_{\mathrm{tail}}=\sum w(\lambda_\chi)$ over these characters, with far weight $S$. The far budget of
the leaf is
\[
  F=\bigl\lceil S(1+\eta)V\bigr\rceil-\1_{\mathrm{inside}}\,n\bigl\lfloor Sw(b)\bigr\rfloor .
\]
Here the last term is present for an inside leaf, that is when $\rho_1\in R_B$. Then
$|\gamma_1|\le C_B/\LL\le1$, so $\rho_1$ and $\overline{\rho_1}$ are the height-one representatives of $\chi_1$ and
$\bchi_1$, and they contribute $n\,w(\lambda_1)\ge n\,w(b)$ to \eqref{eq:farbudget}. A reserved second family is charged either through its own columns, which carry its far
weight, or by subtracting the fixed amount $n_2\lfloor Sw(\hi_2)\rfloor$ from $F$ (\cref{subsec:leafmodel}).
